\section{Results}\label{sec:results}

\subsection{Retrospective Evaluation Setup}

\paragraph{Corpus and labels.}
The released generator (seed 20261029) produced 100 synthetic proposals and 253 planted defects over 23 verifiable rules: 98 absent, 37 negated, 52 placeholder and 66 keyword-present-but-empty (hollow). Thirteen documents have no planted defect; 87 have at least one. Expected verdicts are derived from the planted defects using the frozen policy: 30 Pass, 47 Revise and 23 Fail. Seventeen of the 30 policy-Pass documents contain permissible minor defects, so ``policy Pass'' must not be used as a synonym for ``defect free''. Thirty generated proposals were shaped on the verdict distribution of 30 historical expert opinions; these are not the original private documents or independent expert labels for the generated text.

The corpus and phrase pools were developed with the package. Its self-check guards selected execution-consistency conditions, including unexpected verifier failures on unplanted requirements, but does not require every planted defect to produce a verifier failure. Planted defects can remain pending or be missed. A generator-associated correction to the DOC-002 annex-number capture illustrates its value for development; it also means the corpus is developmental regression evidence, not an untouched semantic-review test set. The released labels are retained for this retrospective analysis, including the lifecycle and format ambiguities discussed in Section~\ref{sec:verifiers}.

\paragraph{Configurations.}
B0 uses rule-specific substrings without the structural and negation handling of the combined system, applying the same aggregation policy to its outputs. S-full uses the verifiers, candidate routes, gate and policy, with a deterministic \emph{flag judge}: it proposes a finding when a candidate carries a negation, assertion-only or placeholder cue. It does not interpret design adequacy. S$-$verifier widens the hybrid statement search to the whole document. B0 does not invoke this judge; the comparison changes multiple mechanisms and cannot isolate a causal typing effect. Memory is present in the software but these outputs do not measure a treatment effect of recalled feedback. No live LLM configuration was run for the reported corpus scores.

\paragraph{Metrics and analysis provenance.}
Alignment is by (document, rule). We recomputed micro-precision, micro-recall and micro-F1 from summed true positives (TP), false positives (FP) and false negatives (FN), and mean per-document F1 from the stored outputs. An empty truth and empty prediction receives per-document F1 of one. We also recomputed policy-verdict accuracy and unweighted Cohen's $\kappa$~\cite{cohen1960,landis1977}, gate acceptance, and referral of missed planted defects to pending review. Out-of-package additions and soft-rule items do not enter the 23-rule detection scores. Published deterministic repetitions have identical core outputs and do not provide independent samples. This revision reaggregates records; it neither regenerates the corpus nor reruns the engine. Previously reported bootstrap intervals and significance tests are not reproduced here and are not retained as verified inference. The document is the observational unit; neither 253 defects nor 2300 rule--document pairs are independent document samples.

\begin{table}[!ht]
\caption{Descriptive results from stored outputs on the fixed developmental corpus. No model capability, isolated typing effect or repaired approval contract is measured.}
\label{tab:results}
\centering\footnotesize
\begin{tabular}{@{}lrrrrrrr@{}}
\toprule
Configuration & TP & FP & FN & Precision & Recall & Micro-F1 & Accuracy / $\kappa$ \\
\midrule
B0 keyword & 47 & 124 & 206 & 0.275 & 0.186 & 0.222 & 28\% / $-0.043$ \\
S-full & 168 & 0 & 85 & 1.000 & 0.664 & 0.798 & 85\% / 0.763 \\
S$-$verifier & 168 & 0 & 85 & 1.000 & 0.664 & 0.798 & 85\% / 0.763 \\
\bottomrule
\end{tabular}
\end{table}

\subsection{Combined-Configuration Trade-offs (RQ2)}

Micro-F1 rises from 0.2217 to 0.7981, a difference of 0.5764. Mean per-document F1 rises from 0.2161 to 0.7798, a paired mean difference of 0.5636. The earlier rounded 0.564 is therefore a different aggregation, not the micro-F1 difference. S-full has perfect per-document F1 on 42 documents and median F1 0.857. The improvement supports this combined configuration over this simple keyword baseline on shared developmental templates. It does not establish superiority over richer deterministic baselines, arbitrary prompted models or human review.

\begin{table}[!ht]
\caption{Retrospective detection by operational tier and severity. Grouping identifies where errors occur; it does not establish that the tier caused the result.}
\label{tab:tiers-results}
\centering\footnotesize
\begin{tabular}{@{}lrrrrrr@{}}
\toprule
& \multicolumn{3}{c}{S-full} & \multicolumn{3}{c}{B0} \\
\cmidrule(lr){2-4}\cmidrule(l){5-7}
Group & TP / FP / FN & Precision & Recall & TP / FP / FN & Precision & Recall \\
\midrule
Hard (10 rules) & 125 / 0 / 17 & 1.000 & 0.880 & 37 / 16 / 105 & 0.698 & 0.261 \\
Hybrid (13 rules) & 43 / 0 / 68 & 1.000 & 0.387 & 10 / 108 / 101 & 0.085 & 0.090 \\
\midrule
Blocker & 17 / 0 / 18 & 1.000 & 0.486 & 3 / 56 / 32 & 0.051 & 0.086 \\
Major & 99 / 0 / 42 & 1.000 & 0.702 & 22 / 53 / 119 & 0.293 & 0.156 \\
Minor & 52 / 0 / 25 & 1.000 & 0.675 & 22 / 15 / 55 & 0.595 & 0.286 \\
\bottomrule
\end{tabular}
\end{table}

Zero false positives is an observed property of these labels and this judge, not a consequence of located evidence. Eighteen of 35 blocker defects are missed. All blocker rules happen to be hybrid in this package; neither that composition nor the recall ordering proves a general severity--verifiability relationship. Seventeen misses occur in hard rules, showing that automatic eligibility can coexist with inadequate proxy predicates.

\subsection{Traceability and Pending Workload (RQ3)}\label{sec:misses}

Sixty-eight of 85 missed planted defects occur in pending hybrid document--rule pairs (80\% referral coverage). Seventeen remain unsurfaced by the evaluated automatic routes. Referral is not recovery: the stored result does not show that a reviewer read the item, identified the defect or corrected the verdict. The full pending set contains 932 unique hybrid pairs, mean 9.32 per document, median nine and range 6--13, excluding soft-rule items. Sixty-eight of these 932 pairs match planted defects (7.3\% label density). The other pairs must not be called false positives: a pending item asserts uncertainty, not a defect. Nor is that density a measure of review time, which was not measured.

\begin{table}[!ht]
\caption{Routes of planted defects. ``Unsufaced'' refers to the evaluated machine routes, not to what a human could find by reading the document. Pending is not a confirmed finding.}
\label{tab:kinds}
\centering\footnotesize
\begin{tabular}{@{}lrrrr@{}}
\toprule
Kind (planted count) & Found & Pending & Unsurfaced & Recall \\
\midrule
Absent (98) & 71 & 27 & 0 & 0.724 \\
Negated (37) & 31 & 6 & 0 & 0.838 \\
Placeholder (52) & 42 & 5 & 5 & 0.808 \\
Hollow (66) & 24 & 30 & 12 & 0.364 \\
\midrule
Total (253) & 168 & 68 & 17 & 0.664 \\
\bottomrule
\end{tabular}
\end{table}

The flag judge submits 53 hybrid findings. Forty-two pass the gate; 11 are rejected, all for \texttt{QUOTE\_TOO\_SHORT}. Gate acceptance is 42/53 (79.2\%). No absence, unlocated or out-of-candidate rejection occurs in these submissions, whose quotes originate in verifier-provided spans. That absence of other rejections does not validate resistance to misleading or fabricated model evidence. The 43 hybrid true positives comprise 42 gate-accepted submissions and one direct SEC-002 structural failure on \texttt{syn-rand-46}; describing all 43 as model judgements would misstate their provenance. Rejected true findings remain pending. Revising the quote floor or overlap predicate would require affected boundary tests and reevaluation, not simply transferring this result to a new gate.

\subsection{Policy-Verdict Errors and Sensitivity}\label{sec:threshold}

S-full and S$-$verifier are identical in these outputs. Widening the region therefore does not expose a difference for this cue-driven judge; it cannot establish the effect of scoping a generative judge. Table~\ref{tab:confusion} shows the S-full recommendations against the policy-derived labels.

\begin{table}[!ht]
\caption{S-full verdict confusion matrix; rows are expected and columns suggested. All 15 disagreements are under-calls.}
\label{tab:confusion}
\centering\small
\begin{tabular}{@{}lrrr@{}}
\toprule
Expected / suggested & Pass & Revise & Fail \\
\midrule
Pass & 30 & 0 & 0 \\
Revise & 6 & 41 & 0 \\
Fail & 4 & 5 & 14 \\
\bottomrule
\end{tabular}
\end{table}

Fifteen of 70 policy-non-Pass documents are under-called (21.4\%); ten receive Pass (14.3\%), comprising six expected Revise and four expected Fail. Of 23 expected Fail documents, four receive Pass (17.4\%) and nine receive a non-Fail recommendation (39.1\%). Ten of the 40 suggested Pass recommendations are incorrect against the policy labels (25\%). The 13 genuinely zero-defect documents all have no finding. Seventeen further documents are policy Pass despite minor planted defects. These denominators prevent both inflating ``clean-document'' performance and mistaking all defects on a Pass document for a false clearance.

With expected labels held fixed under the shipped $\theta=3$ policy, changing the recommendation's minor threshold gives Table~\ref{tab:theta}. The plateau is conditional on this corpus and finding set, not general robustness or a reason to disregard missed blockers.

\begin{table}[!ht]
\caption{Recommendation-threshold sensitivity against the fixed original labels, not against labels regenerated under each threshold.}
\label{tab:theta}
\centering\footnotesize
\begin{tabular}{@{}lrrrrr@{}}
\toprule
$\theta$ & 1 & 2 & 3 (shipped) & 4 & 5 \\
\midrule
Accuracy & 72\% & 82\% & 85\% & 85\% & 85\% \\
Cohen's $\kappa$ & 0.534 & 0.713 & 0.763 & 0.763 & 0.763 \\
\bottomrule
\end{tabular}
\end{table}

\subsection{A Historical Walkthrough and Its Label Boundaries}

The generated proposal \texttt{syn-hist-30} concerns a balance-query capability opened to third-party channels. Its three gate-confirmed blockers are placeholder concurrency control (CAP-001), a negated binding-relation check (SEC-002), and placeholder exposure-risk assessment (SEC-004). A fourth confirmed finding is the hard diagram rule DOC-001, while MON-001 remains pending. The stored expected and suggested verdicts are both Fail, and per-document F1 is 0.889. The automated hybrid actor in this corpus case is the flag judge, not an LLM.

Two parts of this case cannot be interpreted as independent evidence of business defects. DOC-001 fails an image-only diagram under the frozen Mermaid predicate, without assessing what the image shows. MON-001 treats ``New page tracking points must be added to count usage'' as lacking an implementation statement. For a proposal before construction, that may instead be an implementation plan whose completeness needs review; lack of completed deployment is not automatically a proposal defect. We retain these historical labels to keep the analysis version-consistent, but identify the task/label discrepancies. Domain-owner adjudication and a new frozen evaluation are required before treating a corrected predicate as validated.

\subsection{Development-Team Workflow Evidence (RQ4)}\label{sec:pilot}

The authors report five runs on three private technical proposals: one revised/reviewed again and one resumption run. Only two reviews reached a final verdict. Across these runs they report 58 findings: 40 verifier-located, 12 model-judged and six operator additions; three of the 12 model-judged submissions were rejected for out-of-candidate quotations. They also report 18 decisions (15 accept, three dismiss, no rewrite), statistics for 13 rules and 21 exemplars. These categories describe the authors' workflow account, not a reconstructed private-store audit or an independent expert study. The private source documents, model/prompt configuration and full approval/decision trace are not available for independent verification here.

The account supports feasibility of recording and aggregating feedback in ordinary operator-led use. It does not establish verified consent for every record, guaranteed pending disposition, reduced workload, independent reviewer acceptance or a learning curve. It also does not repair the callable-interface and configuration-recovery counterexamples described in Section~\ref{sec:contract}. An anonymous trace tying inputs, gate outcomes, dispositions and effective configuration to each final opinion would be needed to strengthen this workflow evidence without publishing confidential content.

